\documentclass[12pt]{article}
\usepackage[utf8]{inputenc}
\usepackage[T1]{fontenc}
\usepackage{hyperref}
\title{Salawati Language Resources for K-12 Curriculum Development}
\author{Compiled by Claude Sonnet 5.5 (claude-sonnet-5-5) for Andr\'as Kornai}
\date{}
\begin{document}
\maketitle

\section{Introduction}
Salawati (ISO 639-3 xmx; Glottocode made1253), also known as Maden, is an Austronesian language of Salawati Island in the Raja Ampat Islands, Indonesia. Sources cite about 600 speakers (2001) and a status of vulnerable to endangered; younger people mostly have passive knowledge. Dialects named in the sources are Butlih (Banlol, Fiawat), Kawit and Tepin, and it is unclear whether they form one language or two.

\section{Primary Educational Resources}
No primers, textbooks or school programmes for Salawati were found in search results.

\section{Linguistic and Reference Materials}
\begin{itemize}
\item Wikipedia, ``Salawati language'': \url{https://en.wikipedia.org/wiki/Salawati_language}
\item Glottolog map page: \url{https://glottolog.org/resource/languoid/id/made1253.bigmap.html}
\item Language Documentation \&amp; Description journal article that turned up in the search (I did not verify its content or whether it concerns Salawati): \url{https://lddjournal.org/article/1221/galley/2464/download}
\end{itemize}

\section{Religious and Community Resources}
No scripture or community organisation was found.

\section{Digital and Audio Resources}
\begin{itemize}
\item Edinburgh DataShare, ``Data from several Austronesian languages of Raja Ampat (As, Batta, Biga, Ma'ya, Salawati)'': \url{https://datashare.ed.ac.uk/handle/10283/8573}
\end{itemize}

\section{Curriculum Development Context}
Salawati is spoken as a first language by adults only and is not known to be taught in schools. The Edinburgh recordings are the main documentation I found. A curriculum would need to start from them and from new community work, with attention to dialect differences.
\end{document}